% Lecture notes: Week 3, Lecture 2 -- Mechanical Properties
% To hide this lecture from the website, add a line containing only:  % draft
% To set the date shown on the website, add a line like:  % posted: 2026-09-02
\documentclass[11pt]{article}
\usepackage{mech2030notes}
\begin{document}

% ##################################################################
%  WEEK 3, LECTURE 2
% ##################################################################
\lecture{3}{2}{Mechanical Properties}

\noindent
\begin{minipage}[c]{0.55\textwidth}
\centering
\begin{tikzpicture}[scale=0.85]
  \dogbone{0}{0.35}{1.0}
  \node[pin] at (0,0.8) {}; \node[pin] at (0,-0.8) {};
  \draw[dim] (0.9,0.8) -- (0.9,-0.8) node[midway, right] {$\ell_0$};
  \node at (0,0) {\small$\sim\!A_0$};
  \draw[force] (0,2.55) -- (0,3.3) node[above] {$P$};
  \draw[force] (0,-2.55) -- (0,-3.3) node[below] {$P$};
  \node[align=center, font=\small] at (-1.7,0) {initial\\config.};
  \dogbone{3.6}{0.25}{1.5}
  \node[pin] at (3.6,1.2) {}; \node[pin] at (3.6,-1.2) {};
  \draw[dim] (4.5,1.2) -- (4.5,-1.2) node[midway, right] {$\ell$};
  \node[align=center, font=\small] at (3.6,3.6) {deformed config.};
\end{tikzpicture}
\end{minipage}%
\hfill
\begin{minipage}[c]{0.42\textwidth}
\begin{itemize}
  \item $\boxed{\sigma = \dfrac{P}{A}}$
  \item $\boxed{\epsilon = \dfrac{\Delta\ell}{\ell_0} = \dfrac{\ell-\ell_0}{\ell_0}}$
  \item These are ``\emph{engineering}'' stress and strain.
  \item Normalized by the \emph{initial} area $A_0$ and length $\ell_0$.
\end{itemize}
\end{minipage}

\section*{Engineering Stress--Strain Curve}

\noindent
\begin{minipage}[c]{0.55\textwidth}
\centering
\begin{tikzpicture}
  \draw[axis] (0,0) -- (6.4,0) node[right] {$\epsilon$};
  \draw[axis] (0,0) -- (0,3.6) node[above] {$\sigma$};
  \draw[plotline] (0,0) -- (0.8,2.0) .. controls (1.6,2.8) and (2.6,3.0) .. (3.4,3.0)
                  .. controls (4.2,3.0) and (4.8,2.7) .. (5.3,2.4);
  \node at (5.3,2.4) {\Large$\times$};
  \node[right] at (5.45,2.4) {fracture};
  \draw[dashed] (0,2.0) -- (0.8,2.0);
  \draw[dashed] (0,3.0) -- (3.4,3.0);
  \node[left] at (0,2.0) {$\sigma_y$};
  \node[left] at (0,3.0) {$\sigma_u$};
  \node[pin] at (0.8,2.0) {}; \node[pin] at (3.4,3.0) {};
  \draw[thin] (0.2,0.5) -- (0.45,0.5) -- (0.45,1.125);
  \node[right] at (0.45,0.75) {\small$E$};
  \draw[thinarrow] (1.3,2.15) to[bend left=15] (2.4,2.75);
  \node[align=center, font=\small] at (2.3,1.9) {strain\\hardening};
\end{tikzpicture}
\end{minipage}%
\hfill
\begin{minipage}[c]{0.42\textwidth}
\begin{itemize}
  \item The \emph{linear elastic} region has slope
        \[ \boxed{E = \sigma/\epsilon}, \]
        where $E$ is \emph{Young's modulus}.
  \item Plastic yield first occurs at the \emph{yield stress} $\sigma_y$.
  \item The maximum load occurs at the \emph{ultimate stress} $\sigma_u$.
\end{itemize}
\end{minipage}

\section*{Unloading}

\noindent
\begin{minipage}[c]{0.55\textwidth}
\centering
\begin{tikzpicture}
  \draw[axis] (0,0) -- (4.6,0) node[right] {$\epsilon$};
  \draw[axis] (0,0) -- (0,3.6) node[above] {$\sigma$};
  \draw[plotline] (0,0) -- (1.0,2.0) .. controls (1.8,2.6) and (2.6,2.85) .. (3.4,3.0);
  \draw[plotline] (3.4,3.0) -- (1.9,0);
  \draw[thin, dotted] (3.4,3.0) -- (3.4,0);
  \node[pin] at (1.0,2.0) {}; \node[pin] at (3.4,3.0) {};
  \node[above left] at (1.0,2.0) {$\sigma_y$};
  \draw[thinarrow] (0.25,0.9) -- (0.65,1.7);
  \draw[thinarrow] (1.8,2.75) -- (2.8,3.15);
  \draw[thinarrow] (2.9,2.3) -- (2.4,1.3);
  \draw[thin] (0.2,0.4) -- (0.5,0.4) -- (0.5,1.0);
  \node[right] at (0.5,0.65) {\small$E$};
  \draw[thin] (2.1,0.4) -- (2.4,0.4) -- (2.4,1.0);
  \node[right] at (2.4,0.65) {\small$E$};
  \draw[dim] (0,-0.4) -- (3.4,-0.4) node[midway, fill=white] {$\epsilon$};
  \draw[dim] (0,-0.9) -- (1.9,-0.9) node[midway, fill=white] {$\epsilon^p$};
  \draw[dim] (1.9,-0.9) -- (3.4,-0.9) node[midway, fill=white] {$\epsilon^e$};
\end{tikzpicture}
\end{minipage}%
\hfill
\begin{minipage}[c]{0.42\textwidth}
\begin{itemize}
  \item Unloading occurs \emph{elastically}, with slope $E$.
  \item The \emph{total strain} is additively decomposed into elastic and plastic parts:
        \[ \boxed{\epsilon = \epsilon^e + \epsilon^p} \]
\end{itemize}
\end{minipage}

\section*{0.2\% Offset Yield Stress}

\noindent
\begin{minipage}[c]{0.55\textwidth}
\centering
\begin{tikzpicture}
  \draw[axis] (0,0) -- (5.4,0) node[right] {$\epsilon$};
  \draw[axis] (0,0) -- (0,3.6) node[above] {$\sigma$};
  \draw[plotline, name path=curve] plot[domain=0:5, samples=60] (\x,{3.2*(1-exp(-0.9*\x))});
  \draw[thick, name path=off] (0.8,0) -- (2.0,3.456);
  \path[name intersections={of=curve and off, by=Y}];
  \node[pin] at (Y) {};
  \node[below right] at (Y) {$\sigma_y$};
  \draw[thin] (0.8,0.08) -- (0.8,-0.08);
  \node[below, align=center, font=\small] at (0.8,-0.1) {$\epsilon = 0.002$\\(0.2\%)};
  \draw[thin] (0.12,0.2) -- (0.32,0.2) -- (0.32,0.776);
  \node[right] at (0.32,0.45) {\small$E$};
  \draw[thin] (0.95,0.2) -- (1.15,0.2) -- (1.15,0.776);
  \node[right] at (1.15,0.45) {\small$E$};
\end{tikzpicture}
\end{minipage}%
\hfill
\begin{minipage}[c]{0.42\textwidth}
\begin{itemize}
  \item Sometimes the exact $\sigma_y$ is hard to determine.
  \item Take the intercept of a line with slope $E$ passing through $\epsilon = 0.2\%$.
  \item This is the ``\emph{0.2\% offset yield stress}.''
\end{itemize}
\end{minipage}

\section*{Ductile vs.\ Brittle}

\noindent
\begin{minipage}[c]{0.55\textwidth}
\centering
\begin{tikzpicture}
  \draw[axis] (0,0) -- (5.4,0) node[right] {$\epsilon$};
  \draw[axis] (0,0) -- (0,3.4) node[above] {$\sigma$};
  \draw[plotline] (0,0) -- (0.9,2.8);
  \node at (0.9,2.8) {\Large$\times$};
  \node[right] at (1.0,2.9) {Brittle};
  \draw[plotline, red!60!black] (0,0) -- (0.7,1.3) .. controls (1.8,1.9) and (3.5,1.95) .. (4.6,1.75);
  \node at (4.6,1.75) {\Large$\times$};
  \node at (2.6,1.55) {Ductile};
  \draw[thin, dotted] (0.9,2.8) -- (0.9,0);
  \draw[thin, dotted] (4.6,1.75) -- (4.6,0);
  \draw[dim] (0,-0.4) -- (0.9,-0.4) node[midway, below=2pt] {$\epsilon_{f1}$};
  \draw[dim] (0,-1.1) -- (4.6,-1.1) node[midway, below=2pt] {$\epsilon_{f2}$};
\end{tikzpicture}
\end{minipage}%
\hfill
\begin{minipage}[c]{0.42\textwidth}
\begin{itemize}
  \item \emph{Ductile} materials can sustain more plastic strain $\epsilon^p$ before fracture.
  \item They are typically \emph{tougher} than brittle materials, absorbing more energy
        \[ \boxed{\mathcal{E} = \int_0^{\epsilon_f} \sigma(\epsilon)\, d\epsilon} \]
        before fracture.
\end{itemize}
\end{minipage}

\end{document}
